\section{Introduction}
\label{sec:intro}

\subsection{Motivation: bridging physical site evidence and digital truth}

AEC coordination moves constantly between two incompatible views of the same building. On site,
evidence is egocentric and fragmented: a phone photograph, a marked floorplan crop, a short chat
message from a front-line worker. In the project record, truth is allocentric and structured: a BIM
model, IFC element types, topology, and GUIDs. The traceability gap is the missing translation between
these two views. A coordinator does not merely need to know that an image contains \emph{a window};
they need to know which of the many near-identical windows in the BIM cloud should receive the issue,
handoff record, or BCF link. Before an AI system can answer higher-level AEC questions --- what
happened, who is responsible, what action should follow --- it must solve the foundational grounding
question: \emph{where exactly is it in the model?}

The setting is deliberately cold-start. A construction or facilities team has a complete IFC/BIM
model but no history of labelled site imagery; on day one they want to point a phone at an element,
add a short note, and be told \emph{which} modelled element it is. The final step is a selection from
a retrieved candidate pool in which the ground truth is present (median 76, in-pool 100\%), so the
difficulty is not retrieval but \textbf{discrimination among visually identical siblings}, in a regime
with no paired supervision to learn that discrimination from. The natural baseline --- fine-tune a
multimodal reranker to score the pool end-to-end --- is exactly the configuration that plateaus: it
reaches Top-1 6.7\% with the answer already in the pool, because the signal that separates two
identical windows is not in either window's pixels but in the building's relational structure, which a
black-box matcher must reconstruct internally and unreliably (we develop this design rationale in
\secref{sec:rationale}~\citep{sutton2019bitter}).

\subsection{Research questions}

This paper makes one contribution --- an \textbf{ontology-grounded, topology-derived spatial
address} and the calibrated routing that makes it usable. Ontology supplies the floor and the
guardrail: the vocabulary, storey, class, and schema-compliant execution boundary that prevent
hallucinated element IDs. Topology supplies the discriminator: the local relational structure that
separates visually identical siblings. The contribution is developed through three research
questions:

\begin{itemize}
  \item \textbf{RQ1 --- Representation (\secref{sec:eval-oracle}).} \emph{What is the minimal sufficient
  spatial address?} We show it is \textbf{type-conditional}: a coarse ontological prefix (storey +
  class) that is necessary but saturated, completed by a class-specific topological body --- the
  position-slot $(i, M)$ for fillers, a connectivity fingerprint for walls. At an oracle level it
  reorders the pool from Top-1 4.9\% to \textbf{78.5\%}, with the gain partitioned by element type
  (fillers 91\%, walls 64\%).

  \item \textbf{RQ2 --- Mechanism (\secref{sec:eval-realized}).} \emph{How is a noisily-recovered address
  used without losing recall, and how much of the ceiling is realizable?} Hard filtering destroys recall,
  so the address is a \textbf{soft prior in a recall-fixed pool}; one real extractor lifts the
  addressable Top-1 from 6.6\% to \textbf{67.6\%}, its confidence passes an ECE gate (AUROC 0.80),
  and --- since reweighting a finest-grained prior is a no-op --- its payoff is \textbf{selective
  prediction}: deferring the least-confident ${\sim}20\%$ raises answered-set Top-1 to \textbf{80.6\%}
  \citep{guo2017calibration, geifman2017selective}.

  \item \textbf{RQ3 --- Architecture (\secref{sec:eval-depth}).} \emph{Where should the relational context
  live --- extracted deep at inference, or compiled into the node?} A \textbf{depth law}: deeper
  context is informationally more discriminative but its recovery collapses with depth (per-hop
  reliability $0.40 \to 0.05 \to 0$), so the realizable confusable set saturates at one hop (median
  $13 \to 8.2 \to 8.1$). The architecture therefore compiles depth into the node and extracts at
  depth $\le 1$, placing learning at the neural$\to$symbolic interface \citep{mao2019nscl}.
\end{itemize}

\subsection{Related work and positioning}
\label{sec:related}

Prior work approaches the image-to-building-model problem from two directions that do not meet in
the middle.

\paragraph{Vision-language spatial perception.} One line endows VLMs with spatial reasoning by
training on synthetic spatial-QA data; SpatialVLM demonstrates that metric and relational spatial
judgments are learnable from generated supervision \citep{chen2024spatialvlm}, and VLM-based
scene-graph extraction recovers object-relation structure from images.
% TODO [CITATION NEEDED]: the org note also names "Li et al. 2024" and "Wang et al. 2024"
% (pixels-to-graphs VLMs) — could not verify which papers these are; confirm before adding.
These systems answer in \emph{image} space --- a region, a relation, a distance --- and stop at the
image boundary: none of them names an element in a building's authoritative model.

\paragraph{LLM-driven graph retrieval.} A second line connects language models to graph stores.
Document-level GraphRAG builds its graph by LLM extraction from text \citep{edge2024graphrag}, which
is noisy in AEC because element attributes and spatial relationships are rarely verbalized;
text-to-Cypher generation hallucinates node and relationship names and requires LLM-supervised
verification scaffolding to approach reliability \citep{tiwari2024autocypher, gusarov2025graphrag}.
IFC-native systems avoid the noisy-graph problem by building on the schema directly: IFC-graph
converts the model into a queryable property graph \citep{zhu2023ifcgraph}, Graph-RAG pipelines
extract information from IFC data \citep{iranmanesh2025ifcgraphrag}, and MCP4IFC drives IFC-based
design through LLM tool calls \citep{nithyanantham2025mcp4ifc}; Text2BIM generates building models
from language \citep{du2024text2bim}. All of these, however, operate on \emph{structured or textual}
queries --- none grounds unstructured multimodal site evidence to an element.

\paragraph{Positioning.} To our knowledge, this is the first work to combine open-vocabulary
multimodal perception with IFC-native deterministic graph retrieval for \emph{element-level}
grounding in construction site workflows. Prior IFC-based systems
\citep{nithyanantham2025mcp4ifc, iranmanesh2025ifcgraphrag} operate on structured queries; prior
VLM-based spatial systems \citep{chen2024spatialvlm} lack IFC-grounded retrieval. This work bridges
the gap, placing the system in the high-determinism, multimodal-input quadrant of the landscape:
flexible vision-language perception on the input side, ontology-constrained graph execution on the
output side, with a calibrated routing layer between them. The neuro-symbolic precedent --- neural
perception feeding a symbolic executor, with learning at the interface --- is the concept-learner
line \citep{mao2019nscl}.

\subsection{Contributions}

\begin{enumerate}
  \item \textbf{A neuro-symbolic interpreter architecture} that maps unstructured multimodal AEC
  evidence (photo + note + plan) to IFC element GUIDs via a structured intermediate representation
  --- a per-field \texttt{\{value, confidence, source\}} contract --- achieving hallucination-resistant
  retrieval: the neural layer describes evidence, but only the symbolic layer names GUIDs
  (\secref{sec:design}).

  \item \textbf{A type-conditional, image-recoverable spatial address} --- the ordinal position-slot
  for openings, a connectivity fingerprint for walls --- computable from the raw IFC model with no
  human labels, that reorders the retrieved pool from Top-1 4.9\% to 78.5\% at an oracle level and
  underlies the realized 67.6\% (\secref{sec:eval-oracle}, \secref{sec:eval-realized}).

  \item \textbf{An oracle-based evaluation methodology} that decomposes retrieval failure into the
  symbolic-layer ceiling versus neural extraction quality --- demonstrating 100\% ground-truth
  retention, isolating extraction as the bottleneck, and yielding a recall-safety principle (union,
  never intersection, when combining mature and immature signals) and a measured depth law for
  relational context (\secref{sec:eval}).

  \item \textbf{An empirical analysis of where learned and deterministic components belong}: a
  per-field profile of a LoRA fine-tuned VLM showing saturation on coarse fields and failure on
  discriminating ones, motivating delegation to deterministic visual specialists with calibrated,
  selectively-predictive routing (ECE gate, defer-on-low-confidence) rather than further fine-tuning
  (\secref{sec:eval-neural}, \secref{sec:eval-realized}).
\end{enumerate}

\subsection{Scope, stated up front}

Two honesty boundaries frame every result. First, the diagnostic ceilings (RQ1, and the oracle rows
of the evaluation) assume perfect extraction; the realizable numbers come from one deterministic
extractor (the position-slot), with the wall fingerprint and remaining descriptors demonstrated at
oracle level only --- full realization is future work. Second, all measurements are on a single
synthetic project (cold-start by design); the \emph{form} of the contribution --- a type-conditional,
shallow, image-recoverable, calibrated-soft address --- is general, but the specific reliabilities are
project-specific. To keep the claim anchored beyond our own ablations, we add standard lexical and
dense retrieval baselines and a triage-effort proxy: generic retrievers plateau at Top-1 1.7\% /
Top-10 15.6--25.0\%, the realized position-slot specialist closes the filler subset to 67.6\% Top-1,
and the perfect-address ceiling reduces expected inspection effort from a median 38 candidates to 0.5.
We do not claim to beat an end-to-end matcher in the large-data limit; we claim that in the cold-start
regime we can actually measure, the structured address is what makes grounding realizable, auditable,
and transferable.
